% =============================================================================
% Springer Publication Figure Environments (LNCS `llncs.cls` / Journal `sn-jnl.cls`)
% Auto-generated by `scripts/generate-springer-svgs.clj`
% Tip: Use `\usepackage{svg}` with `\includesvg[width=\textwidth]{...}` or convert
%      each SVG to PDF via `inkscape --export-type=pdf` and use `\includegraphics`.
% =============================================================================

\begin{figure}[t]
  \centering
  \includegraphics[width=\textwidth]{figures/springer/fig1_architecture_and_dsl.pdf}
  \caption{\textbf{System architecture of the typed executable optimizer laboratory.}
  \textbf{(a)} Binder-free, sort-stratified Domain-Specific Language (DSL) over six syntactic roles
  (\texttt{:operation}, \texttt{:predicate}, \texttt{:score}, \texttt{:count}, \texttt{:initializer},
  and \texttt{:operator}), guaranteeing 100\% valid-by-construction execution under typed generation,
  subtree mutation, and homologous crossover.
  \textbf{(b)} Unchanged shared core verified via locked SHA-256 hashes across domains (Hypothesis H5),
  enforcing a dual resource ledger: an outer 256-assessment-slot allowance (where duplicates and
  fallbacks consume slots) and an inner 256-proposal / 4,096-step runtime cap (where cache hits still
  charge one proposal unit).
  \textbf{(c--d)} Domain adapters for exact-$K$ cardinality-constrained portfolio selection
  ($N=50, K=5, T=365$) and 30-job single-machine permutation scheduling.}
  \label{fig:architecture}
\end{figure}

\begin{figure}[t]
  \centering
  \includegraphics[width=\textwidth]{figures/springer/fig2_pareto_fronts_2panel.pdf}
  \caption{\textbf{Bi-objective Pareto fronts, dominated decision clouds, and hypervolume geometry.}
  \textbf{(a)} Exact-$K$ portfolio selection ($N=15, K=5$, all $\binom{15}{5}=3{,}003$ subsets) minimizing
  negative mean daily return ($-\mu_p$) and daily return standard deviation ($\sigma_p$).
  \textbf{(b)} Single-machine scheduling (30 jobs, $B=256$ proposals) minimizing mean completion time
  and total tardiness alongside classical dispatching heuristics (SPT: blue diamond; EDD: purple square).
  In both panels, initial random solutions are marked by amber triangles, dominated solutions by grey dots,
  and nondominated Pareto-optimal solutions by emerald circles.
  Because both decision spaces ($\binom{N}{K}$ subsets and $n!$ permutations) are finite combinatorial sets,
  the true Pareto set $P^\star$ is a finite collection of discrete points rather than a continuous manifold.
  We therefore display two complementary geometric objects: (i) a \emph{smooth Fritsch--Carlson monotone cubic
  Hermite curve} (solid dark-emerald line) illustrating the continuous shape-preserving trade-off trend without
  artificial overshoot, and (ii) the \emph{exact 2-D Hypervolume attainment region} (light emerald shading bounded
  by a dashed orthogonal staircase), which reflects the axis-aligned dominance cones $\bigcup_{p \in A}[p, r]$
  integrated with respect to the reference point $r$.}
  \label{fig:pareto-fronts}
\end{figure}

\begin{figure}[t]
  \centering
  \includegraphics[width=\textwidth]{figures/springer/fig3_five_arms_and_prior_ablation.pdf}
  \caption{\textbf{Held-out test normalized hypervolume across all five outer search arms and generator priors.}
  Comparison of evolved executable optimizers (\texttt{A1: GP-Mixed}, \texttt{A2: GP-Mutation-Only}),
  random program sampling (\texttt{A3: Random-DSL}), fixed-scaffold expression evolution
  (\texttt{A4: Formula-Only}), and configured classical multi-objective baselines
  (\texttt{A5: Configured-Conventional}, NSGA-II / PLS) on \textbf{(a)} Portfolio Selection ($N=50, K=5, T=365$)
  and \textbf{(b)} Single-Machine Scheduling (30 jobs).
  Grouped bars compare the naive uniform prior (\texttt{:development-only-v1}, amber with triangle markers)
  against the capacity-preserving prior (\texttt{:budget-aware-v2}, emerald with circle markers), which
  eliminates outer-slot starvation and achieves practically superior test hypervolume gains ($p < 0.0001$).}
  \label{fig:five-arms-and-prior}
\end{figure}

\begin{figure}[t]
  \centering
  \includegraphics[width=\textwidth]{figures/springer/fig4_anytime_budget_scaling.pdf}
  \caption{\textbf{Anytime performance scaling across inner solution proposal budgets $B \in \{16, 32, 64, 128, 256\}$.}
  Frozen winning optimizers and configured baselines evaluated on held-out test instances at five logarithmic
  proposal budgets on \textbf{(a)} Portfolio Selection and \textbf{(b)} Single-Machine Scheduling.
  Distinct geometric marker shapes (circle, square, diamond, triangle, cross) and stroke dash styles ensure
  legibility in both color and grayscale print.}
  \label{fig:budget-scaling}
\end{figure}

\begin{figure}[t]
  \centering
  \includegraphics[width=\textwidth]{figures/springer/fig5_mechanism_ablations_3panel.pdf}
  \caption{\textbf{Secondary mechanism ablations verifying the sources of evolutionary and runtime performance.}
  \textbf{(a)} Per-attempt syntactic/sort validation hit rate (\%) for typed sort-directed variation
  (\texttt{:budget-aware-v2}, emerald) versus untyped grammar-blind tree variation with post-hoc validation
  (\texttt{:untyped-validated}, vermilion), where untyped generation and mutation suffer $>90\%$ rejection rates.
  \textbf{(b)} Surgical AST knockouts (\texttt{K0}--\texttt{K6}) on frozen winning portfolio optimizers:
  deleting the evolved search body (\texttt{K1}) collapses hypervolume by $-0.491$, blinding the score expression
  (\texttt{K3}) drops hypervolume by $-0.129$, and collapsing bi-objective trade-offs to a single objective
  (\texttt{K6}) drops hypervolume by $-0.037$, whereas neutralizing the initializer (\texttt{K2}) changes
  hypervolume by only $-0.001$.
  \textbf{(c)} Inner candidate memoization cache-off (\texttt{:cache? false}, grey) versus cache-on
  (\texttt{:cache? true}, emerald): because cache hits still charge one proposal unit, held-out test
  hypervolume is bitwise identical ($|\Delta\text{HV}| = 0.0$) while physical objective evaluations are
  reduced by 14.5\%--16.5\%.}
  \label{fig:mechanism-ablations}
\end{figure}

\begin{figure}[t]
  \centering
  \includegraphics[width=\textwidth]{figures/springer/fig6_meta_evolution_and_ast_structure.pdf}
  \caption{\textbf{Outer meta-evolution trajectory and syntactic structure of the 30 final evolved programs.}
  \textbf{(a)} Mean best-so-far training normalized hypervolume across the 8 outer assessment batches
  ($8 \times 32 = 256$ slots, $N=30$ independent seeds) for \texttt{A1: GP-Mixed}, \texttt{A2: GP-Mutation-Only},
  and \texttt{A4: Formula-Only}, demonstrating substantial multi-generation improvement over batch 0
  ($+0.047$ to $+0.049$ training hypervolume).
  \textbf{(b)} Frequency tally of syntactic roles, operators, and features across the 30 final selected
  \texttt{GP-Mixed} programs on Portfolio Selection (emerald, 30/30 syntactically unique) and Single-Machine
  Scheduling (purple, 30/30 syntactically unique), showing automatic domain-appropriate operator selection
  (\texttt{:swap} and \texttt{:restart} in Portfolio versus \texttt{:insert}, \texttt{:swap}, and \texttt{:first}
  in Scheduling) alongside bi-objective (\texttt{:objective-0} + \texttt{:objective-1}) and non-stationary
  progress-adaptive (\texttt{:progress}) scoring formulas.}
  \label{fig:meta-evolution-and-ast}
\end{figure}
